\documentclass[10pt,a4paper]{amsart}
\usepackage[margin=19mm]{geometry}
\usepackage{amsmath,amssymb,mathtools}
\usepackage[scr=rsfs,cal=euler]{mathalfa}
\usepackage{xfrac}
\usepackage[unicode,hidelinks,bookmarks=false]{hyperref}
\newcommand{\ie}{i.e.\ }
\newcommand{\cf}{cf.\ }
\newcommand{\resp}{respectively\ }
\newcommand{\Subsection}{\subsection}
\providecommand{\nobreakdash}{\nobreak\hbox{-}}
\setlength{\emergencystretch}{2em}
\multlinegap0pt
\usepackage{amssymb}
\usepackage{mathtools}
\usepackage{tikz}
\usepackage{float}
\usepackage[shortlabels]{enumitem}
\newtheorem{definition}{Definition}
\newtheorem{proposition}{Proposition}
\newtheorem{notation}{Notation}
\newtheorem{lemma}{Lemma}
\usepackage{cleveref}
\crefname{definition}{Definition}{Definitions}
\crefname{proposition}{Proposition}{Propositions}
\crefname{notation}{Notation}{Notations}
\crefname{lemma}{Lemma}{Lemmas}
\newcommand{\la}{\lambda}
\DeclareMathOperator{\p}{\mathcal{P}}
\title{Enumeration of Odd-Dimensional Partitions Modulo 4}
\author{Aditya Khanna}
\date{Source: arXiv:2207.07513v4 (CC BY 4.0)}
\begin{document}
\maketitle

\noindent\emph{Source and licence.} Enumeration of Odd-Dimensional Partitions Modulo 4. Version arXiv:2207.07513v4 (CC BY 4.0), distributed by arXiv under the Creative Commons Attribution 4.0 International licence, http://creativecommons.org/licenses/by/4.0/, as recorded in the arXiv licence metadata for this exact version and shown on the abstract page https://arxiv.org/abs/2207.07513v4. Full licence notice: \texttt{licenses/arxiv-2207.07513-CC-BY-4.0.txt}.\par\smallskip

\noindent\textit{Editorial scope.} Counted unit: the article's lemma lem:8 (source0.tex lines 1114--1122). The article's complete original proof of it is delivered with it (source0.tex lines 1123--1178, one \texttt{proof} environment of 56 lines). No mathematics is written, repaired or re-derived: the statement and the proof are the article's own lines, in the article's own notation, and no retained line is substituted or re-flowed.  Retained uncounted as the source-local context the proof needs: lines 492--494; lines 559--564; lines 763--765.  Nothing was omitted from any retained span, so the package carries no omission record.  No retained line contains a citation, so the wrapper carries no bibliography and no bibliography entry.  No retained line contains a figure environment, an image-inclusion command or drawing code, so no image is delivered and the package declares no asset dependency.  Editorial formatting only, in addition: an AMS \texttt{amsart} preamble that restates the article's own declarations and macros used by the retained lines, namely the environments \texttt{definition}, \texttt{proposition}, \texttt{notation}, \texttt{lemma} on separate counters, together with the standard packages those lines need.  The retained environments print in this document as Definition 1, Proposition 1, Notation 1, Lemma 1 in order of appearance, whereas the article prints the same items with the numbering of its own full document.  The complete native source of the article as distributed in the arXiv e-print for this exact version is bundled unchanged as \texttt{sources/128282/source0.tex}.
\begin{definition}\label{notn:aff}
For partitions $\mu$ and $\la$ such that $\la$ is the $2^R$-parent of $\mu$, define the \textit{affected hook-length} $h^\la_\mu$ to be the element of $H(\la)$ such that replacing $h^\la_\mu$ by $h^\la_\mu - 2^R$ yields a $\beta$-set of $\mu$. In our rim hook removal notation, $H(\mu) \sim_\beta H(\la)[h^\la_\mu \to h^\la_\mu - 2^R]$.
\end{definition}

\begin{proposition}\label{prop:inversions}
Let $\lambda$ be a $2^R$-parent of $\mu$ with $|\mu| < 2^R$.  Then
\[
\eta^{\lambda}_{\mu} = N_{\lambda}(h^\lambda_\mu) - \mathbb{I}_{H(\lambda)}(h^\lambda_\mu - 2^{R-1}) +  \mathbb{I}_{H(\lambda)}(h^\lambda_\mu + 2^{R-1}) + \mathbb{I}_{H(\lambda)}(h^\lambda_\mu - 2^R-2^{R-1}). 
\]
\end{proposition}

\begin{notation}\label{notn:valid-values}
For $r \in \{1,\ldots, 2^R\}$, define $\lambda_{[r]}$ to be the partition for which $H(\la_{[r]})$ is constructed by replacing 0 by $2^R$ in $H(\mu)^{+r}$. If $2^R\in H(\mu)^{+r}$, define $\lambda_{[r]} =\mu$. A value of $r$ is called \textit{valid} if $\lambda_{[r]} \neq \mu$.
\end{notation}

\begin{lemma}\label{lem:8}
Let $\mu\vdash 2^{R-1}$ be an odd-dimensional partition of $2^{R-1}$ with $H(\mu) = \{2^{R-1},b, b-1,\ldots, 1\}$ for $R\geq 2$ and $b\geq 0$. Let $\p_2(\mu)$ denote the set of Type II $2^R$-parents of $\mu$, then 
\[
S_{\p_2(\mu)}({\mu}) = \begin{cases}
2 ,& \text{if }\ell(\mu)\text{ is even}\\
3, & \text{if }\ell(\mu) \text{ is odd}.\\
\end{cases}
\]
\end{lemma}

\begin{proof}
Denote by $\lambda_{[r]}$ the partition for which $H(\la_{[r]})$ is obtained by replacing 0 with $2^R$ in $ H(\mu)^{+r}$ when $2^R\not\in H(\mu)^{+r}$. $\p_2(\mu)$ is the set of such $\la_{[r]}$. For an invalid shift $s$ (\cref{notn:valid-values}), we let $\la_{[s]}$ be $\mu$. First, for all $\lambda\in\p_2(\mu)$, the affected hook-length (\cref{notn:aff}), $h^\lambda_\mu$, is  $2^R$, and so $s_2(h^\lambda_\mu) = 1$. This reduces the computation of the signed sum as follows:
\[
S_{\p_2(\mu)}({\mu}) = \sum\limits_{\lambda\in \p_2(\mu)} (-1)^{\eta^\lambda_\mu + 1} = (-1)\cdot\sum\limits_{\lambda\in \p_2(\mu)} (-1)^{\eta^\lambda_\mu}.
\]
As $2^R - (2^R+ 2^{R-1}) < 0$,   \cref{prop:inversions} tells us that $
\eta^\lambda_\mu = N_{\lambda}(2^R) - \mathbb{I}_{H(\lambda)}(2^{R-1}) + \mathbb{I}_{H(\lambda)}(2^R +2^{R-1})$ for any $\la\in \p_2(\mu)$.

In order to deduce whether $2^{R-1}$, $2^R$, or $2^{R}+2^{R-1}$ occur in $H(\la)$, we divide the set $\{1,2,\ldots, 2^R\}$ of values that $r$ can take into six intervals, of which some may be empty. As one may check, the following cases work for all values of $b\geq 0$.
\begin{enumerate}
\item Consider the case when {$1 \leq r\leq 2^{R-1} - b - 1$}. The largest entry in $H(\mu)^{+r}$ is $2^{R-1} + r$ which is at most $ 2^{R} - b -1$ and is thus strictly smaller than $2^R$. The second largest element in $H(\mu)^{+r}$ must be $b+r$ which is strictly smaller than $2^{R-1}$. Therefore, the maximum element of $H(\la_{[r]})$ is $2^R$ and $H(\la_{[r]})$ does not contain $2^{R-1}$. This forces $\mathbb{I}_{H(\lambda)}(2^{R-1})$ and $ \mathbb{I}_{H(\lambda)}(2^R +2^{R-1})$ to be zero. So, we have $\eta^{\lambda_{[r]}}_\mu = N_{\lambda_{[r]}}(2^R)$ which is the number of entries smaller than $2^R$. This is exactly one less than the number of entries in $H(\la)$. The length of $\la$ is $r$ more than the length of $\mu$ which gives us $\eta^{\lambda_{[r]}}_\mu = \ell(\mu) + r - 1$.

\item Consider {$2^{R-1} - b \leq r \leq 2^{R-1}-1$}. If $b = 0$, then the interval in this case is empty. If $b$ is positive, then $H(\mu)^{+r}$ contains $2^{R-1}$ which gives $\mathbb{I}_{H(\lambda)}(2^{R-1})=1$. On the other hand, as the largest value in $H(\mu)^{+r}$ is $b+r$, which is strictly smaller than $2^R + 2^{R-1}$, $\mathbb{I}_{H(\lambda)}(2^R + 2^{R-1})=0$. So $\eta^{\lambda_{[r]}}_\mu = N_{\lambda_{[r]}}(2^R) - 1$. We find $N_{\lambda_{[r]}}(2^R) \ell(\mu) + r-1$ as in part 1 and obtain $\eta^{\lambda_{[r]}}_\mu = \ell(\mu) + r - 2$.
\item When $r = 2^{R-1}$, we get $2^{R-1} + 2^{R-1} = 2^R\in H(\mu)^{+2^{R-1}}$, so $\lambda_{[2^{R-1}]} \not\in \p_2(\mu)$.

\item In the case when $2^{R-1} + 1 \leq r \leq 2^{R}-b-1$, $H(\lambda_{[r]})$ always contains $2^{R-1}$, so $\mathbb{I}_{H(\lambda)}(2^{R-1}) = 1$. Furthermore, the largest element of $H(\la_{[r]})$ is $2^{R-1} + r\leq 2^{R-1} + 2^R - b -1$ which is strictly smaller than $2^R + 2^{R-1}$, and thus $\mathbb{I}_{H(\lambda)}(2^R +2^{R-1}) = 0$. The number of elements in $H(\la_{[r]})$ are $\ell(\mu) + r$, out of which one element is $2^R$ and the other element is $2^{R-1} + 2^{R} - b - 1$, which is greater than $2^R$. This shows that the number of entries in $H(\la_{[r]})$ strictly smaller than $2^R$, that is $N_{\la_{[r]}}(2^R)$, is $\ell(\mu) + r - 2$. From this, we obtain $\eta^{\lambda_{[r]}}_\mu = (\ell(\mu)+r -2) - 1 = \ell(\mu) + r - 3$.

\item When $2^R - b \leq r \leq 2^{R}-1$, $2^{R}\in H(\mu)^{+r}$, so $\lambda_{[r]}\not\in \p_2(\mu)$.
\item Finally, when $r = 2^R$, the set $H(\mu)^{+2^R}$ contains both $2^{R-1}$ and $2^R + 2^{R-1}$. Furthermore, the set contains $1,2,\ldots, 2^{R}-1$, which means $N_{\la_{[2^R]}}(2^R) = 2^R - 1$.
So, $\eta_{\mu}^{\lambda_{[2^R]}} = N_{\lambda_{[2^R]}}(2^R) = (2^{R} - 1) -1 + 1 = 2^R -1$.
\end{enumerate}

We sum over the intervals in (1), (2), (4) and (6) to give
\begin{align*}
S_{\p_2(\mu)}({\mu}) &= (-1)\cdot\Bigg(\sum\limits_{r = 1}^{2^{R-1}-b-1} (-1)^{\ell(\mu)+r-1} + \sum\limits_{r = 2^{R-1}-b}^{2^{R-1}-1} (-1)^{\ell(\mu)+r-2} \\ 
&\qquad\qquad\quad+ \sum\limits_{r=2^{R-1} +1}^{2^{R}-b-1} (-1)^{\ell(\mu)+r - 3} + (-1)^{2^R - 1}\Bigg).
\end{align*}
We multiply the $-1$ inside the sum to find
\begin{align*}
S_{\p_2(\mu)}({\mu}) &= \sum\limits_{r = 1}^{2^{R-1}-b-1} (-1)^{\ell(\mu)+r} + \sum\limits_{r = 2^{R-1}-b}^{2^{R-1}-1} (-1)^{\ell(\mu)+r-1} \\ 
&\qquad\qquad\quad+ \sum\limits_{r=2^{R-1} +1}^{2^{R}-b-1} (-1)^{\ell(\mu)+r - 2} + (-1)^{2^R}.
\end{align*}
The last term is equal to $+1$. From the rest of the terms, we factor out $(-1)^{\ell(\mu)}$ to get

\[
 S_{\p_2(\mu)}({\mu}) = (-1)^{\ell(\mu)} \Bigg(\sum\limits_{r = 1}^{2^{R-1}-b-1} (-1)^{r} + \sum\limits_{r = 2^{R-1}-b}^{2^{R-1}-1} (-1)^{r-1} + \sum\limits_{r=2^{R-1} +1}^{2^{R}-b-1} (-1)^{r-2}\Bigg) + 1.
\]
The exponent $r-2$ in the third sum can be replaced with $r$ while keeping the same bounds. This yields
\[
 S_{\p_2(\mu)}({\mu}) = (-1)^{\ell(\mu)} \Bigg(\sum\limits_{r = 1}^{2^{R-1}-b-1} (-1)^{r} + \sum\limits_{r = 2^{R-1}-b}^{2^{R-1}-1} (-1)^{r-1} + \sum\limits_{r=2^{R-1} +1}^{2^{R}-b-1} (-1)^{r}\Bigg) + 1,
\]
and we notice that the third sum is the same as the first sum but the bounds are shifted up by $2^{R-1}$. Adding $2^{R-1}$ to the exponent of $(-1)^r$ does not change the sign of the output, so the first and third sums are equal. We now have 
\[
S_{\p_2(\mu)}({\mu})  = 1 + 2(-1)^{\ell(\mu)}\sum\limits_{r = 1}^{2^{R-1}-b - 1} (-1)^{r} + \sum\limits_{r = 2^{R-1}-b}^{2^{R-1}-1} (-1)^{r-1}.
\]
The $2^{R-1}$ in the bounds of the sum on the right hand side does not effect the value of the sum as $R\geq 2$. We reconsider this sum as ranging from $-b$ to $-1$, and the exponent of $-1$ ranges from $-(b+1)$ to $-2$. As $-1$ is its own inverse, we may write the sum as ranging over $x' = 2$ to $b+1$. A shift of 2 preserves the parity, and thus we may rewrite the sum as ranging over $x = 0$ to $x = b-1$.
We have 
\[
S_{\p_2(\mu)}({\mu})  = 1 + 2(-1)^{\ell(\mu)}\sum\limits_{r = 1}^{2^{R-1}-b - 1} (-1)^{r} + \sum\limits_{x=0}^{b-1} (-1)^{x}.
\]
The number of elements in $H(\mu)$ is one more than $b$, and so $b=\ell(\mu)-1$. This allows us to rewrite our sums in terms of $\ell(\mu)$ as
\[
S_{\p_2(\mu)}({\mu})  = 1 + 2(-1)^{\ell(\mu)}\sum\limits_{r = 1}^{2^{R-1}-\ell(\mu)} (-1)^{r} + \sum\limits_{x=0}^{\ell(\mu) - 2} (-1)^x.
\]
When $\ell(\mu)$ is even, we get $S_{\p_2(\mu)}({\mu}) = 1 + 2\cdot 1\cdot 0 + 1 = 2$. When $\ell(\mu)$ is odd, we get $S_{\p_2(\mu)}({\mu})  = 1 + 2\cdot (-1) \cdot (-1) + 0 = 3$.
\end{proof}

\end{document}
